
One then obtains the following result ([Per75], II, Th.~\SGARefLink{prop:sga3-VIA-2.4}{2.4}; see also
[Per76], Prop.~4.1.1).

\medskip
\phantomsection\label{thm:sga3-VIA-2.6.5}%
\noindent\textbf{Theorem 2.6.5.}
\textit{Let \(k\) be a field, and let \(G\) be a \(k\)-group scheme.}
\begin{enumerate}
\item[\SGAInlineRefTarget{thm:sga3-VIA-2.6.5:item:i}(i)] \textit{There is a unique group subscheme \(G^0\) of
\(G\), called the \emph{identity component} of \(G\), such that:}
  \begin{enumerate}
  \item[\SGAInlineRefTarget{thm:sga3-VIA-2.6.5:item:a}(a)] \textit{the underlying space of \(G^0\) is the
  irreducible component of the identity element;}
  \item[\SGAInlineRefTarget{thm:sga3-VIA-2.6.5:item:b}(b)] \textit{\(G^0\to G\) is a flat closed immersion; that is,
  \[
    \mathcal O_{G,g}=\mathcal O_{G^0,g}
    \qquad\text{for every }g\in G^0.
  \]}
  \end{enumerate}
\item[\SGAInlineRefTarget{thm:sga3-VIA-2.6.5:item:ii}(ii)] \textit{Moreover, \(G^0\) is quasi-compact,
geometrically irreducible, and is a characteristic subgroup of \(G\).}
\item[\SGAInlineRefTarget{thm:sga3-VIA-2.6.5:item:iii}(iii)] \textit{If \(G\) is connected, then \(G=G^0\).}
\end{enumerate}

\noindent\textit{Proof.}
\SGARefLink{thm:sga3-VIA-2.6.5:item:i}{(i)} Recall first that \(G\) is separated
(\hyperref[sga3:VIA:0.3]{0.3}), and hence, a fortiori,
quasi-separated. Let \(U\) be an affine open subset of \(G\) containing
the generic point \(\omega\) of \(C^0\). By
\hyperref[rem:sga3-VIA-2.6.3]{2.6.3} and
\hyperref[prop:sga3-VIA-2.6.4]{2.6.4},
\[
  \overline U=C^0U
\]
is both open and closed, and, set-theoretically,
\[
  C^0=\overline{\{\omega\}}
\]
is the intersection of these open-and-closed subsets.

As \(U\) runs through the affine open subsets containing \(\omega\),
the \(G\)-schemes \(\overline U\) form a projective system whose
transition morphisms are affine, since they are closed immersions.
We can therefore form its projective limit \(G^0\)
(cf.\ EGA~IV\textsubscript{3}, 8.2.2). Thus, for every affine open
\(V\) of \(G\), \(G^0\cap V\) is the spectrum of the algebra
\[
  \varinjlim_U\mathcal O_G(V\cap\overline U)
  =
  \mathcal O_G(V)\Big/\sum_U I_{\overline U}(V),
\]
where \(I_{\overline U}(V)\) denotes the kernel of
\[
  \mathcal O_G(V)\longrightarrow
  \mathcal O_G(V\cap\overline U).
\]
It follows that the underlying space of \(G^0\) is \(C^0\), and that
\(G^0\to G\) is a closed immersion.

Moreover, for every \(g\in G^0\), \(\mathcal O_{G^0,g}\) is the
inductive limit, as \(V\) runs through the affine open subsets of
\(G\) containing \(g\), of the \(k\)-algebras
\[
  \mathcal O_{G^0}(V\cap G^0)
  =
  \varinjlim_U\mathcal O_G(V\cap\overline U).
\]
This double inductive limit is identified with
\[
  \varinjlim_U\
  \varinjlim_{\substack{V\\g\in V\subset\overline U}}
  \mathcal O_G(V)
  =
  \mathcal O_{G,g}.
\]
Thus
\[
  \mathcal O_{G^0,g}=\mathcal O_{G,g}
\]
(see also EGA~IV\textsubscript{2}, 5.13.3(ii)), and consequently
\(i:G^0\to G\) is a flat closed immersion. Conversely, this condition
implies that \(i^*(\mathcal O_G)=\mathcal O_{G^0}\), so \(G^0\) is
uniquely determined by conditions \SGARefLink{thm:sga3-VIA-2.6.5:item:a}{(a)} and \SGARefLink{thm:sga3-VIA-2.6.5:item:b}{(b)}. This proves \SGARefLink{thm:sga3-VIA-2.6.5:item:i}{(i)}.

The first two assertions of \SGARefLink{thm:sga3-VIA-2.6.5:item:ii}{(ii)} follow from
\hyperref[cor:sga3-VIA-2.6.2]{2.6.2}. Finally, let \(S\) be a
\(k\)-scheme and let \(\phi\) be an automorphism of the \(S\)-group
\(G_S\). For every \(s\in S\), \(\phi_s\) sends
\(G^0_{\kappa(s)}\) into itself, and therefore
\(\phi(G^0_S)\subset G^0_S\). Moreover, the closed immersion
\[
  i_S:G^0_S\hookrightarrow G_S
\]
obtained from \(i\) by base change is flat, so
\[
  \mathcal O_{G_S,\phi(z)}
  =
  \mathcal O_{G^0_S,\phi(z)}
  \qquad\text{for every }z\in G^0_S.
\]
Hence \(\phi\circ i_S\) factors through \(G^0_S\). This proves that
\(G^0\) is a characteristic subgroup of \(G\), and proves \SGARefLink{thm:sga3-VIA-2.6.5:item:ii}{(ii)}.
Finally, \SGARefLink{thm:sga3-VIA-2.6.5:item:iii}{(iii)} is a special case of part \SGARefLink{thm:sga3-VIA-2.6.5:item:i}{(i)} of the following
proposition.%

\medskip
\phantomsection\label{prop:sga3-VIA-2.6.6}%
\noindent\textbf{Proposition 2.6.6.}
\textit{Let \(k\) be a field, and let \(G\) be a \(k\)-group acting
on a \(k\)-scheme \(X\) in such a way that
\[
  G\times X\longrightarrow X\times X,\qquad
  (g,x)\longmapsto(gx,x)
\]
is surjective. Suppose that \(X\) is quasi-separated. Then:}
\begin{enumerate}
\item[\SGAInlineRefTarget{prop:sga3-VIA-2.6.6:item:i}(i)] \textit{Every connected component \(C\) of \(X\) is
irreducible.}
\item[\SGAInlineRefTarget{prop:sga3-VIA-2.6.6:item:ii}(ii)] \textit{Let \(\eta\) be the generic point of \(C\), and
let \(L\) be the algebraic closure of \(k\) in \(\kappa(\eta)\).
Then \(C\) is an \(L\)-scheme, is geometrically irreducible over
\(L\), and the morphism
\[
  G^0_L\times_LC\longrightarrow C\times_LC
\]
is surjective.}
\item[\SGAInlineRefTarget{prop:sga3-VIA-2.6.6:item:iii}(iii)] \textit{In particular, if \(C\) contains a rational
point \(x\), then \(C\) is geometrically irreducible over \(k\), and
\[
  \phi:G^0\longrightarrow C,\qquad g\longmapsto gx,
\]
is surjective.}
\end{enumerate}

\noindent\textit{Proof.}
\SGARefLink{prop:sga3-VIA-2.6.6:item:i}{(i)} Let \(C\) be a connected component of \(X\), let \(Y\) be an
irreducible component of \(X\) contained in \(C\), let \(\eta\) be
the generic point of \(Y\), and let \(U\) be an affine open subset of
\(X\) containing \(\eta\). Since \(X\) is quasi-separated, \(U\) is
retrocompact in \(X\). Hence, by
\hyperref[prop:sga3-VIA-2.6.4]{2.6.4},
\[
  \overline U=G^0U
\]
is an open-and-closed subset of \(X\) which meets \(C\), and therefore
contains \(C\). By \hyperref[rem:sga3-VIA-2.6.3]{2.6.3}, the
intersection of the \(\overline U\), as \(U\) runs through a
fundamental system of affine open neighborhoods of \(\eta\), is
\(\overline{\{\eta\}}\). It follows that
\[
  C=\overline{\{\eta\}},
\]
which proves \SGARefLink{prop:sga3-VIA-2.6.6:item:i}{(i)}.

The first assertion of \SGARefLink{prop:sga3-VIA-2.6.6:item:ii}{(ii)}, and also that of \SGARefLink{prop:sga3-VIA-2.6.6:item:iii}{(iii)}, follows from
\hyperref[prop:sga3-VIA-2.6.1]{2.6.1}. We first prove the second
assertion of \SGARefLink{prop:sga3-VIA-2.6.1:item:iii}{(iii)}. Let \(\omega\), respectively \(\eta\), denote the
generic point of \(G^0\), respectively \(C\). Let \(z\in C\), and set
\(K=\kappa(z)\). Since \(G^0\), respectively \(C\), is geometrically
irreducible over \(k\), the scheme \(G^0_K\), respectively \(C_K\),
is irreducible; denote its generic point by \(\xi\), respectively
\(\beta\). We saw in the proof of
\hyperref[prop:sga3-VIA-2.6.4]{2.6.4} that
\[
  \mu_z:G_K\longrightarrow C_K,\qquad h\longmapsto h\cdot z,
\]
sends \(\xi\) to \(\beta\). Likewise, \(\phi(\omega)=\eta\).

Put \(L=\kappa(\xi)\), and let \(\xi_L\) and \(z_L\) be the \(L\)-points
obtained from \(\xi\) and \(z\). Then
\[
  \xi_L\cdot z_L=\beta'
\]
is a point of \(C_L\) lying above \(\beta\in C_K\), and hence also
above \(\eta\in C\). Consider the cartesian square
\begin{center}
  \phantomsection\label{diag:sga3-VIA-2.6.6-cartesian-square}%
  \begin{tikzcd}[row sep=large,column sep=large]
    G_L^0 \arrow[r,"\varphi_L"] \arrow[d,"\pi_{G^0}"'] &
    C_L \arrow[d,"\pi_C"] \\
    G^0 \arrow[r,"\varphi"'] &
    C
  \end{tikzcd}
\end{center}
Since
\[
  \phi_L\bigl(\pi_{G^0}^{-1}(\omega)\bigr)
  =
  \pi_C^{-1}(\eta)
\]
(cf.\ EGA~I, 3.4.8), there is a \(g\in G^0_L\) such that
\(\phi_L(g)=\beta'\). Therefore
\[
  z_L
  =
  \xi_L^{-1}\cdot\phi_L(g)
  =
  \phi_L(\xi_L^{-1}g),
\]
whence
\[
  \phi\bigl(\pi_{G^0}(\xi_L^{-1}g)\bigr)
  =
  \pi_C(z_L)
  =
  z.
\]
This proves that \(\phi\) is surjective.

We now prove the second assertion of \SGARefLink{prop:sga3-VIA-2.6.6:item:ii}{(ii)}. It is enough to show that,
for every \(z\in C\), the morphism
\[
  \mu_z:
  G^0_L\otimes_L\kappa(z)
  \longrightarrow
  C\otimes_L\kappa(z)
\]
is surjective. But this follows from \SGARefLink{prop:sga3-VIA-2.6.6:item:iii}{(iii)}, since \(z\) is a rational
point of \(C\otimes_L\kappa(z)\).

\medskip
\phantomsection\label{rem:sga3-VIA-2.6.7}%
\noindent\textbf{Remark 2.6.7.---}
\textit{Under the hypotheses of \hyperref[prop:sga3-VIA-2.6.6]{2.6.6}, if \(C(k)=\varnothing\), the
morphism
\[
  G^0\times_kC\longrightarrow C\times_kC
\]
need not be surjective. For example, let \(k=\mathbb R\) and
\(G=\{\pm1\}_{\mathbb R}\). The \(G\)-torsor
\[
  X=\operatorname{Spec}\mathbb R[X]/(X^2+1)
\]
is connected, but
\[
  G^0\times_{\mathbb R}X\longrightarrow X\times_{\mathbb R}X
\]
is not surjective. (Here \(L=\mathbb C\), and
\[
  G^0\times_{\mathbb R}X\longrightarrow X\times_{\mathbb C}X
\]
is an isomorphism.)}
